El marco metodológico se constituye como el resultado de la aplicación sistemática y lógica de los conceptos y fundamentos abordados en el marco teórico. Azuero et al. \cite{azuero_2019} lo definen como un conjunto de acciones destinadas a describir y analizar el fondo del problema planteado, mediante procedimientos específicos que incluyen técnicas de observación y recolección de datos, los cuales determinarán el cómo se realizará la investigación. 

Este capítulo se dividirá en secciones que definan la estructura del marco metodológico, con el fin de sustentar los métodos y recursos elegidos para el desarrollo de la investigación, asegurando que cada fase esté alineada con los principios teóricos fundamentales y que los resultados obtenidos sean fiables, replicables y efectivos en el tratamiento de la problemática planteada.